\section{The pure conductor-four character at the quarter residues}
\label{sec:pure-conductor4-quartic-strict}

\paragraph{Scope.}
This argument treats the original finite character $\chi_{-4}$
with trivial cubic part.  The real quartic class input is inherited
from AC-G152-class.  Its $K$-twisted companion has an elementary
finite Stickelberger certificate, so no new analytic interval,
class-number selector or full fundamental-unit basis is needed.
One existing real-quartic unit supplies both required local images.

Put $p=269$, $K=\mathbb Q(\sqrt{-3})$, and
$\delta=\chi_{-4}|_{G_K}$.
Use the same symbol for the specified rational character when
discussing abelian character fields.
Write $c_p$ for the actual cyclotomic character and $\omega_p$
for its Teichmuller part.  Fix an embedding of $\mathbb Q(i)$
in $\mathbb Q_p$ with $i\equiv82\pmod p$.
For a discrete coefficient $D$, the full strict group is
\[
 \begin{gathered}
 S_{\rm str}(K,D)=\ker\bigl(H^1_{\rm cts}(G_K,D)
                                  \longrightarrow\mathcal L(D)\bigr),\\
 \mathcal L(D)=\prod_{v\mid p}H^1(G_{K_v},D)
              \times\prod_{v\nmid p,\infty}H^1(I_v,D).
 \end{gathered}
 \tag{P4Q.1}
\]

\begin{theorem}[AC-G162: the pure conductor-four quarter-residue rows]
\label{AC-G162}
Let $E/\mathbb Q_p$ be any finite extension and let $O$ be its
original ring of integers.  Let $m\in\mathbb Z$, $e\geq1$, and
let $\kappa:G_K\to1+\mathfrak m_O$ be continuous and unramified
at every finite place away from $p$.
If $e+m\equiv66$ or $200\pmod{268}$, then
\[
 \begin{gathered}
 D=(E/O)(\delta\omega_p^{-m}\kappa^{-1}c_p^{-e}),\\
                         S_{\rm str}(K,D)=0.
 \end{gathered}
 \tag{P4Q.2}
\]
The residual strict groups vanish over every finite field of
characteristic $p$.  Arbitrary ramification of the original
coefficient ring is permitted.
For $m=0,\kappa=1$, this supplies additional weights for the
already present pure finite label $\delta$.
The proof establishes the precise companion class and unit
statements needed for this trivial cubic part.
\end{theorem}

\begin{proof}
\proofstep{1. The real quartic and its exact companion.}
Put $\psi=\omega_p^{67}$, so $\psi(2)=-i$, and let
$\xi=\delta\psi$ as a rational character.
Its cyclic quartic field, with the orientation fixed in
AC-G152-class, is
\[
 \begin{gathered}
 F=\mathbb Q(w),\qquad w^4-pw^2+25p=0,\\
 s=\frac{2w^2-p}{13},\qquad s^2=p,\\
 b(w)=\frac{w^3-219w}{65},\qquad b^2(w)=-w,
                           \qquad\xi(b)=-i.
 \end{gathered}
 \tag{P4Q.3}
\]
Here $F$ is totally real, since $\xi$ is even; its unique
quadratic subfield is $\mathbb Q(s)$.
It is disjoint from $K$, and $N=KF$ has degree eight over
$\mathbb Q$ and degree four over $K$.
Extend $b$ to fix $K$, and let $c$ fix $F$ and conjugate $K$.
These automorphisms commute, with orders four and two.

Write $\eta_K=\chi_{-3}$ and
$\xi'=\eta_K\xi=\chi_{12}\psi$.
It is primitive of conductor $f=12p=3228$ and is odd.
Indeed $\chi_{12}=\chi_{-3}\chi_{-4}$ is primitive of conductor
twelve, and its modulus is coprime to the primitive modulus $p$
of $\psi$, so both local conductor factors remain.
Set $z=\sqrt{-3}\,w$.  Then
\[
 \begin{gathered}
 F'=\mathbb Q(z),\qquad z^4+3pz^2+225p=0,\\
 F'=N^{\langle cb^2\rangle},\qquad
 b(z)=-\frac{z^3+657z}{195},\qquad\xi'(b)=-i.
 \end{gathered}
 \tag{P4Q.4}
\]
The displayed polynomial is Eisenstein at $p$, so $[F':\mathbb Q]=4$.
The element $cb^2$ fixes $z$ and is the exact kernel of the
twist on $C_2\times C_4$; this proves the field assertion.
The CRT representative $1885$ is one modulo twelve and two
modulo $p$.  It fixes $K$ and $i$, and induces the chosen $b$
on both quartic fields.  In particular the extension of $b$ to
$N$ is not assigned using a CRT representative that conjugates $K$.
Both quartics have quadratic subfield $\mathbb Q(s)$.

Let $a$ be the residual coefficient character in (P4Q.2), and
write $r\equiv e+m\pmod{268}$.
The residual Kummer quotient and inverse character are
\[
 \begin{array}{c|c|c}
 r&\eta=a/\bar c_p&\eta^{-1}(b)\\ \hline
 66&\bar\xi^{-1}|_{G_K}&187=-\bar i\\
 200&\bar\xi|_{G_K}&82=\bar i.
 \end{array}
 \tag{P4Q.5}
\]
Thus $N/K$ is the required quartic splitting cover.
The finite character in the original coefficient remains $\delta$;
neither $\psi$ nor $\xi$ is added as a new original finite label.

\proofstep{2. The companion class certificate by an integral annihilator.}
Let $L=F(\sqrt[3]p)$ be the degree-twelve field of AC-G152-class.
That theorem proves $\operatorname{Cl}(L)[p]=0$.
Extension of ideals followed by norm is multiplication by three
from $F$, and by six from $\mathbb Q(s)$.
Both integers are prime to $p$, so
\[
              \operatorname{Cl}(F)[p]=
              \operatorname{Cl}(\mathbb Q(s))[p]=0.
 \tag{P4Q.6}
\]

For the companion, put $C=\mathbb Q(\zeta_f)$ with $f=3228$.
Its Galois group has order $\varphi(f)=1072$, prime to $p$,
and $[C:F']=268$ is also prime to $p$.
For $a$ prime to $f$ let $\sigma_a(\zeta_f)=\zeta_f^a$.
The precise Stickelberger element and integral annihilator are
\[
 \begin{aligned}
 \Theta&=\frac1f\sum_{\substack{1\leq a\leq f\\(a,f)=1}}
                                a\,\sigma_a^{-1},\\
 A&=(5-\sigma_5)\Theta
   =\sum_{\substack{1\leq a\leq f\\(a,f)=1}}
                     \left\lfloor\frac{5a}{f}\right\rfloor\sigma_a^{-1}.
 \end{aligned}
 \tag{P4Q.7}
\]
The second equality follows by reindexing multiplication by five
modulo $f$; it proves integrality before reduction at $p$.
By Stickelberger's theorem $A$ annihilates $\operatorname{Cl}(C)$.
The exact applied source is Sharifi, Definition~4.5.1,
Lemma~4.5.2 and Theorem~4.5.6: the integral intersection
$\mathbb Z[G]\Theta\cap\mathbb Z[G]$ annihilates the class group.
Its sign convention gives $(\sigma_5-5)\Theta$; its negative
in (P4Q.7) lies in the same ideal.

Here is a fully finite evaluation with no division modulo $p$.
The element two generates $\mathbb F_p^\times$.
For $a\not\equiv0\pmod p$, if $a\equiv2^j\pmod p$, then
$\psi(a)=(-i)^j$.
The other factor $\chi_{12}$ has values $1$ on residues $1,11$
modulo twelve, $-1$ on $5,7$, and zero on nonunits.
Integer summation gives
\[
 \begin{aligned}
 \sum_{a=1}^f a\xi'(a)&=6456+45192i,\\
 B_1(\xi')=\frac1f\sum_{a=1}^f a\xi'(a)&=2+14i,
       \qquad \xi'(5)=-1,\\
 \sum_{\substack{1\leq a\leq f\\(a,f)=1}}
       \left\lfloor\frac{5a}{f}\right\rfloor\xi'(a)^{-1}
                                      &=12-84i.
 \end{aligned}
 \tag{P4Q.8}
\]
The last integer sum is checked directly, separately from the
weighted sum.  On the two odd quartic characters $\xi'$ and
$\xi'^{-1}$ the annihilator scalars are respectively
\[
       12-84i\equiv118,\qquad12+84i\equiv175\pmod p.
 \tag{P4Q.9}
\]
They are units.  The $p$-torsion of $\operatorname{Cl}(C)$
in each of these character components is therefore zero.
Extension from $F'$ to $C$ is injective on $p$-torsion, because
norm after extension is $268$; it is also equivariant for the
abelian Galois actions.  Thus the $b^2=-1$ part of
$\operatorname{Cl}(F')[p]$ vanishes.
The $b^2=1$ part descends by the norm-extension operator
$1+b^2$ to $\mathbb Q(s)$, whose $p$-class group is zero by
(P4Q.6).  On that part the operator is multiplication by two.
Consequently
\[
                         \operatorname{Cl}(F')[p]=0.
 \tag{P4Q.10}
\]
This proves the needed companion statement without an analytic
regulator or a proposed value for its class number.

\proofstep{3. The exact class sector of $N$ and Kummer descent.}
Norm to $F=N^{\langle c\rangle}$ followed by extension gives
$1+c=0$ on $\operatorname{Cl}(N)[p]$.
Norm to $F'=N^{\langle cb^2\rangle}$ gives $1+cb^2=0$.
On the $b^2=-1$ part, these operators sum to two.
Hence
\[
                \operatorname{Cl}(N)[p]^{b^2=-1}=0.
 \tag{P4Q.11}
\]
No vanishing of a different class sector of $N$ is inferred.

Put $k=O/\mathfrak m_O$.
Prime-to-$p$ restriction and inflation, for $N/K$ and for every
local decomposition or inertia quotient, identify the residual
strict group with
\[
 \begin{gathered}
 e_{\eta^{-1}}\ker\left(V_{N,S}\otimes_{\mathbb F_p}k
                        \longrightarrow\mathcal P_N\otimes_{\mathbb F_p}k\right),\\
 V_{N,S}=\{x\in N^\times:p\mid v(x)\text{ for }v\nmid p\}/N^{\times p},\\
 \mathcal P_N=\bigoplus_{v\mid p}N_v^\times/N_v^{\times p}.
 \end{gathered}
 \tag{P4Q.12}
\]
Over $N$ the coefficient is $k(1)$; the extra quotient twist
is $\eta$, so invariants select the inverse character.
No adjoining of $\mu_p$ is required.
The full local power condition makes the valuations above $p$
divisible by $p$ as well.  Thus the kernel lies in the all-primes
valuation group with exact sequence
\[
 0\longrightarrow\mathcal O_N^\times/\mathcal O_N^{\times p}
 \longrightarrow V_N\longrightarrow\operatorname{Cl}(N)[p]
 \longrightarrow0.
 \tag{P4Q.13}
\]
Both rows in (P4Q.5) have $b^2=-1$, so (P4Q.11) leaves only
their unit components.

\proofstep{4. One exact global unit and its full local images.}
The quartic polynomial in (P4Q.3) is Eisenstein at $p$.
Together with the unramified quadratic completion of $K$, it
gives the unique $p$-completion of $N$, of indices $(4,2)$.
The element $w$ is a uniformizer, $v_w(p)=4$, and its power
basis over $\mathcal O_{K_p}$ is integral and complete.

Use the existing source unit two of AC-G152-class:
\[
 U=\frac{3464441-655959w-14375w^2+2721w^3}{13}.
 \tag{P4Q.14}
\]
Exact rational multiplication verifies
\[
                   U^4-471078U^3-107707606U^2+471078U+1=0.
 \tag{P4Q.15}
\]
Monicity and constant coefficient one prove that $U$ and its
inverse are integral.  This degree-four relation cubed is the
full degree-twelve relation of the original unit file.
The verifier checks both the rational matrix characteristic
polynomial and its Horner evaluation; no assertion that this
unit generates a complete unit lattice is used.

All denominators in (P4Q.14) are retained before reduction.
In the truncated local ring $\mathbb F_{p^2}[w]/(w^4)$,
the exact values are
\[
 \begin{aligned}
 U&\equiv82+155w+53w^2+230w^3,\\
 \ell=\log\bigl(U/\operatorname{Teich}(\bar U)\bigr)
      &\equiv202w+108w^3,\\
 b(w)&\equiv187w+149w^3.
 \end{aligned}
 \tag{P4Q.16}
\]
The logarithm is computed by solving $U\ell'=U'$ and integrating
through degree three; all required denominators are prime to $p$.
The cubic term of $b(w)$ is retained.
For $\lambda=82,187$, let
$P_\lambda=\sum_{j=0}^3\lambda^{-j}b^j$ on unit classes modulo $p$.
The complete nonlinear substitution gives
\[
                 P_{82}\ell=241w^3,
        \qquad P_{187}\ell=w+191w^3.
 \tag{P4Q.17}
\]
Both values are nonzero and satisfy the indicated eigenvalue
identities.  Integer representatives of the coefficients define
products of the actual global conjugate units before localization.
Changing representatives by multiples of $p$ changes neither
their Kummer classes nor these truncated logarithms.

Since $4<p-1$, logarithm and exponential identify principal
units with the maximal ideal.  The logarithm of a principal
$p$th power has valuation at least five; (P4Q.17) has leading
degree three or one.  Teichmuller units are $p$th powers, and a
root of a principal unit has residue one by bijectivity of
Frobenius on the residue field.
Thus these are nonzero images in the full local Kummer quotient.

\proofstep{5. Unit multiplicities and the original coefficient lift.}
The four archimedean places above a fixed complex embedding of
$K$ form a regular $C_4=\operatorname{Gal}(N/K)$ set.
The equivariant unit character is the regular character minus
the trivial character.  There is no $p$-torsion in the roots of
unity, since $268\nmid[N:\mathbb Q]=8$.
Because $p\nmid4$, these multiplicities persist modulo $p$.
Each of the $b$-eigenvalues $82,187$ has a one-dimensional unit
space over $\mathbb F_p$.
The nonzero images (P4Q.17) make equivariant localization injective
on each space, also after scalar extension to every finite $k$.
Equations (P4Q.11)--(P4Q.13) prove residual strict zero.

Let $\varpi$ be a uniformizer of the original $O$.
The coefficient $D$ of (P4Q.2) is divisible and
$D[\varpi]\simeq k(a)$.
At $p$, $\delta$ is unramified and $\kappa$ reduces to one.
The residual inertia characters $\bar c_p^{-66}$ and
$\bar c_p^{-200}$ have respective orders $134$ and $67$;
both are nontrivial because $K_p/\mathbb Q_p$ is unramified.
An inertia scalar differs from one by an $O$-unit, so
$D^{G_{K_p}}=0$.
Away from $p$, the cyclotomic, Teichmuller and $\kappa$ factors
are unramified, and quadratic inertia invariants are either
zero or the whole divisible module $D$.
Thus $D^J/\varpi D^J=0$ for every local group in (P4Q.1).

The multiplication-by-$\varpi$ long exact sequence gives
local injectivity $H^1(J,D[\varpi])\to H^1(J,D)$.
A strict global class killed by $\varpi$ lifts to a residual
class, and local injectivity makes that lift strict, hence zero.
Every continuous cocycle from compact $G_K$ to discrete $E/O$
has finite image killed by a power of $\varpi$.
Any nonzero strict class would therefore produce a nonzero
one killed by $\varpi$, a contradiction.
This proves (P4Q.2) over the original, arbitrarily ramified ring.
\end{proof}

\paragraph{Actual label, exact count and remaining scope.}
For $m=0,\kappa=1$, the finite label is just $\chi_{-4}|_{G_K}$.
It already factors through $\mathbb Q(i)$ in every prescribed
host with $u\geq2$ and has trivial cubic part.
The new weights are $e=66+134n$, $n\geq0$, so for $1\leq e<T$
their exact number is $\lfloor(T+67)/134\rfloor$.
Together with AC-G160, the union of these proved progressions is
$e\equiv-1\pmod{67}$ for this finite label, with count
$\lfloor T/67\rfloor$.
No positive derived index is asserted zero; the original local
comparison cost and common-window transport remain separate.

\paragraph{Evidence and source scope.}
The complete finite verifier is
\path{verify_pure_conductor4_quartic.py}, with output
\path{pure_conductor4_quartic_certificate.json}.
It verifies the exact integer weighted and floor sums, the two
unit annihilator scalars, both field automorphisms and one full
global-unit relation with both local projectors.
The original AC-G152 unit file and class proof are copied unchanged.
Sharifi's complete Definition~4.5.1 statement, Lemma~4.5.2 proof
and Theorem~4.5.6 proof were freshly read in the retained primary
chapter, with exact applied passages preserved alongside it.
Kummer theory retains the audited Stacks tag~03PK scope; the
finite-cover and rank-one Artin--Tate dictionary retains SGA~1,
Expos\'e~V, \S\S4 and 7.  Equivariant units and class norm/extension
are used at their previously checked scopes.
No new analytic replay, complete unit basis, full class assertion
for $N$, mixed realization, all-prime bound, H04.4 or Lean execution
is claimed.
